
\documentclass{article}
\usepackage{colortbl}
\usepackage{makecell}
\usepackage{multirow}
\usepackage{supertabular}

\begin{document}

\newcounter{utterance}

\centering \large Interaction Transcript for game `function\_detective', experiment `list\_passive\_examples', episode 97 with Qwen3.6{-}35B{-}A3B{-}FP8{-}without{-}reasoning.
\vspace{24pt}

{ \footnotesize  \setcounter{utterance}{1}
\setlength{\tabcolsep}{0pt}
\begin{supertabular}{c@{$\;$}|p{.15\linewidth}@{}p{.15\linewidth}p{.15\linewidth}p{.15\linewidth}p{.15\linewidth}p{.15\linewidth}}
    \# & \multicolumn{2}{c}{Player} && \multicolumn{2}{c}{Game Master} \\
    \hline

    \theutterance \stepcounter{utterance}  
    & & & \multicolumn{4}{p{0.6\linewidth}}{
        \cellcolor[rgb]{0.9,0.9,0.9}{
            \makecell[{{p{\linewidth}}}]{
                \texttt{\tiny{[P1$\langle$GM]}}
                \texttt{I have a hidden Python function for you to figure out.} \\
\\ 
\texttt{Your job is to infer what it does and then submit a correct Python implementation.} \\
\\ 
\texttt{Function signature:} \\
\texttt{(items: List{[}int{]}) {-}> int} \\
\\ 
\texttt{You have 15 turns total.} \\
\\ 
\texttt{I control which valid input{-}output examples you see. You can request them one at a time.} \\
\\ 
\texttt{You must make your best SOLVE attempt no later than the second{-}to{-}last turn.} \\
\texttt{Do not wait until the very end.} \\
\\ 
\texttt{You can use these actions:} \\
\texttt{{-} JSON format described at the end of the prompt} \\
\texttt{{-} SOLVE: ```python ... ```} \\
\\ 
\texttt{Do not add import statements in your guess.} \\
\texttt{Anything you might need is already available, so just write the function body in hypothesis and return the answer.} \\
\\ 
\texttt{Preloaded examples:} \\
\texttt{1. ({[}{-}1, 0, 1{]}, {-}1)} \\
\\ 
\texttt{Scoring:} \\
\texttt{{-} correctness is required for any non{-}zero score} \\
\texttt{{-} in interactive modes, solving with fewer information requests is better} \\
\\ 
\texttt{Input formatting reminders:} \\
\texttt{{-} strings must be quoted} \\
\texttt{{-} booleans must be True/False} \\
\texttt{{-} use commas between multiple values} \\
\texttt{{-} do not write any extra action tags in the same message} \\
\\ 
\texttt{Format example:} \\
\texttt{You write:} \\
\texttt{\{'current\_mode': <probing/confirming/don't know; choose whichever best describes your current state in the task>, 'input\_rationale': <brief description of why you chose this hypothesis>, 'hypothesis': <either the exact string "unknown" if you have no guess yet OR your current best guess as runnable Python matching the given signature, with a one{-}line docstring stating what it does. Never write placeholder code, if you cannot commit to a guess, use "unknown". Match the given signature.} \\
\texttt{e.g.} \\
\texttt{def f(x):} \\
\texttt{	"""Returns the absolute value of x."""} \\
\texttt{	return x if x >= 0 else {-}x} \\
\\ 
\texttt{Use "unknown" if you have no guess.>, 'confidence\_score': <score between 1{-}100. 100 means highest confidence in your hypothesis>\}} \\
\\ 
\texttt{Output only one valid json per turn.} \\
            }
        }
    }
     \\ \\

    \theutterance \stepcounter{utterance}  
    & \multicolumn{4}{p{0.6\linewidth}}{
        \cellcolor[rgb]{0.561,0.737,0.561}{
            \makecell[{{p{\linewidth}}}]{
                \texttt{\tiny{[P1$\rangle$GM]}}
                \texttt{\{'current\_mode': 'probing', 'input\_rationale': 'I have one example: ({[}{-}1, 0, 1{]}, {-}1). This is not enough to determine the function. I need more examples to understand the pattern. Let me request another example.', 'hypothesis': 'unknown', 'confidence\_score': 10\}} \\
            }
        }
    }
    & & \\ \\

    \theutterance \stepcounter{utterance}  
    & & & \multicolumn{4}{p{0.6\linewidth}}{
        \cellcolor[rgb]{0.9,0.9,0.9}{
            \makecell[{{p{\linewidth}}}]{
                \texttt{\tiny{[P1$\langle$GM]}}
                \texttt{Game Violation: ParseError: Response is not valid JSON.} \\
\texttt{Output ONLY:} \\
\texttt{One JSON with confidence\_score, current\_mode, input\_rationale, and hypothesis keys ...} \\
\texttt{OR} \\
\texttt{SOLVE: ```python ... ```} \\
            }
        }
    }
     \\ \\

    \theutterance \stepcounter{utterance}  
    & \multicolumn{4}{p{0.6\linewidth}}{
        \cellcolor[rgb]{0.561,0.737,0.561}{
            \makecell[{{p{\linewidth}}}]{
                \texttt{\tiny{[P1$\rangle$GM]}}
                \texttt{\{"current\_mode": "probing", "input\_rationale": "I have one example: ({[}{-}1, 0, 1{]}, {-}1). This is not enough to determine the function. I need more examples to understand the pattern. Let me request another example.", "hypothesis": "unknown", "confidence\_score": 10\}} \\
            }
        }
    }
    & & \\ \\

    \theutterance \stepcounter{utterance}  
    & & & \multicolumn{4}{p{0.6\linewidth}}{
        \cellcolor[rgb]{0.9,0.9,0.9}{
            \makecell[{{p{\linewidth}}}]{
                \texttt{\tiny{[P1$\langle$GM]}}
                \texttt{OUTPUT: ({[}2, 1{]}, 0)} \\
            }
        }
    }
     \\ \\

    \theutterance \stepcounter{utterance}  
    & \multicolumn{4}{p{0.6\linewidth}}{
        \cellcolor[rgb]{0.561,0.737,0.561}{
            \makecell[{{p{\linewidth}}}]{
                \texttt{\tiny{[P1$\rangle$GM]}}
                \texttt{\{"current\_mode": "probing", "input\_rationale": "I have two examples now: ({[}{-}1, 0, 1{]}, {-}1) and ({[}2, 1{]}, 0). Let me analyze these. For {[}{-}1, 0, 1{]}, the result is {-}1. For {[}2, 1{]}, the result is 0. This doesn't immediately suggest a simple min/max or sum. Let me try another example to gather more data.", "hypothesis": "unknown", "confidence\_score": 20\}} \\
            }
        }
    }
    & & \\ \\

    \theutterance \stepcounter{utterance}  
    & & & \multicolumn{4}{p{0.6\linewidth}}{
        \cellcolor[rgb]{0.9,0.9,0.9}{
            \makecell[{{p{\linewidth}}}]{
                \texttt{\tiny{[P1$\langle$GM]}}
                \texttt{OUTPUT: ({[}3, 2, 1{]}, 0)} \\
            }
        }
    }
     \\ \\

    \theutterance \stepcounter{utterance}  
    & \multicolumn{4}{p{0.6\linewidth}}{
        \cellcolor[rgb]{0.561,0.737,0.561}{
            \makecell[{{p{\linewidth}}}]{
                \texttt{\tiny{[P1$\rangle$GM]}}
                \texttt{\{"current\_mode": "probing", "input\_rationale": "I have three examples now: ({[}{-}1, 0, 1{]}, {-}1), ({[}2, 1{]}, 0), and ({[}3, 2, 1{]}, 0). Let me analyze these. For {[}{-}1, 0, 1{]}, the result is {-}1. For {[}2, 1{]}, the result is 0. For {[}3, 2, 1{]}, the result is 0. This still doesn't immediately suggest a simple min/max or sum. Let me try another example to gather more data.", "hypothesis": "unknown", "confidence\_score": 30\}} \\
            }
        }
    }
    & & \\ \\

    \theutterance \stepcounter{utterance}  
    & & & \multicolumn{4}{p{0.6\linewidth}}{
        \cellcolor[rgb]{0.9,0.9,0.9}{
            \makecell[{{p{\linewidth}}}]{
                \texttt{\tiny{[P1$\langle$GM]}}
                \texttt{OUTPUT: ({[}1, 1{]}, 0)} \\
            }
        }
    }
     \\ \\

    \theutterance \stepcounter{utterance}  
    & \multicolumn{4}{p{0.6\linewidth}}{
        \cellcolor[rgb]{0.561,0.737,0.561}{
            \makecell[{{p{\linewidth}}}]{
                \texttt{\tiny{[P1$\rangle$GM]}}
                \texttt{\{"current\_mode": "probing", "input\_rationale": "I have four examples now: ({[}{-}1, 0, 1{]}, {-}1), ({[}2, 1{]}, 0), ({[}3, 2, 1{]}, 0), and ({[}1, 1{]}, 0). Let me analyze these. For {[}{-}1, 0, 1{]}, the result is {-}1. For {[}2, 1{]}, the result is 0. For {[}3, 2, 1{]}, the result is 0. For {[}1, 1{]}, the result is 0. This still doesn't immediately suggest a simple min/max or sum. Let me try another example to gather more data.", "hypothesis": "unknown", "confidence\_score": 40\}} \\
            }
        }
    }
    & & \\ \\

    \theutterance \stepcounter{utterance}  
    & & & \multicolumn{4}{p{0.6\linewidth}}{
        \cellcolor[rgb]{0.9,0.9,0.9}{
            \makecell[{{p{\linewidth}}}]{
                \texttt{\tiny{[P1$\langle$GM]}}
                \texttt{OUTPUT: ({[}{]}, 0)} \\
            }
        }
    }
     \\ \\

    \theutterance \stepcounter{utterance}  
    & \multicolumn{4}{p{0.6\linewidth}}{
        \cellcolor[rgb]{0.561,0.737,0.561}{
            \makecell[{{p{\linewidth}}}]{
                \texttt{\tiny{[P1$\rangle$GM]}}
                \texttt{\{"current\_mode": "probing", "input\_rationale": "I have five examples now: ({[}{-}1, 0, 1{]}, {-}1), ({[}2, 1{]}, 0), ({[}3, 2, 1{]}, 0), ({[}1, 1{]}, 0), and ({[}{]}, 0). The empty list returns 0. Let me try an example with a single element to see if that helps clarify the pattern.", "hypothesis": "unknown", "confidence\_score": 50\}} \\
            }
        }
    }
    & & \\ \\

    \theutterance \stepcounter{utterance}  
    & & & \multicolumn{4}{p{0.6\linewidth}}{
        \cellcolor[rgb]{0.9,0.9,0.9}{
            \makecell[{{p{\linewidth}}}]{
                \texttt{\tiny{[P1$\langle$GM]}}
                \texttt{OUTPUT: ({[}1, 2, 3{]}, 0)} \\
            }
        }
    }
     \\ \\

    \theutterance \stepcounter{utterance}  
    & \multicolumn{4}{p{0.6\linewidth}}{
        \cellcolor[rgb]{0.561,0.737,0.561}{
            \makecell[{{p{\linewidth}}}]{
                \texttt{\tiny{[P1$\rangle$GM]}}
                \texttt{\{"current\_mode": "probing", "input\_rationale": "I have six examples now: ({[}{-}1, 0, 1{]}, {-}1), ({[}2, 1{]}, 0), ({[}3, 2, 1{]}, 0), ({[}1, 1{]}, 0), ({[}{]}, 0), and ({[}1, 2, 3{]}, 0). Let me analyze these. For {[}{-}1, 0, 1{]}, the result is {-}1. For {[}2, 1{]}, the result is 0. For {[}3, 2, 1{]}, the result is 0. For {[}1, 1{]}, the result is 0. For {[}{]}, the result is 0. For {[}1, 2, 3{]}, the result is 0. This still doesn't immediately suggest a simple min/max or sum. Let me try another example to gather more data.", "hypothesis": "unknown", "confidence\_score": 60\}} \\
            }
        }
    }
    & & \\ \\

    \theutterance \stepcounter{utterance}  
    & & & \multicolumn{4}{p{0.6\linewidth}}{
        \cellcolor[rgb]{0.9,0.9,0.9}{
            \makecell[{{p{\linewidth}}}]{
                \texttt{\tiny{[P1$\langle$GM]}}
                \texttt{OUTPUT: ({[}0{]}, 0)} \\
            }
        }
    }
     \\ \\

    \theutterance \stepcounter{utterance}  
    & \multicolumn{4}{p{0.6\linewidth}}{
        \cellcolor[rgb]{0.561,0.737,0.561}{
            \makecell[{{p{\linewidth}}}]{
                \texttt{\tiny{[P1$\rangle$GM]}}
                \texttt{\{"current\_mode": "probing", "input\_rationale": "I have seven examples now: ({[}{-}1, 0, 1{]}, {-}1), ({[}2, 1{]}, 0), ({[}3, 2, 1{]}, 0), ({[}1, 1{]}, 0), ({[}{]}, 0), ({[}1, 2, 3{]}, 0), and ({[}0{]}, 0). Let me analyze these. For {[}{-}1, 0, 1{]}, the result is {-}1. For {[}2, 1{]}, the result is 0. For {[}3, 2, 1{]}, the result is 0. For {[}1, 1{]}, the result is 0. For {[}{]}, the result is 0. For {[}1, 2, 3{]}, the result is 0. For {[}0{]}, the result is 0. This still doesn't immediately suggest a simple min/max or sum. Let me try another example to gather more data.", "hypothesis": "unknown", "confidence\_score": 70\}} \\
            }
        }
    }
    & & \\ \\

    \theutterance \stepcounter{utterance}  
    & & & \multicolumn{4}{p{0.6\linewidth}}{
        \cellcolor[rgb]{0.9,0.9,0.9}{
            \makecell[{{p{\linewidth}}}]{
                \texttt{\tiny{[P1$\langle$GM]}}
                \texttt{OUTPUT: ({[}{-}4{]}, {-}4)} \\
            }
        }
    }
     \\ \\

    \theutterance \stepcounter{utterance}  
    & \multicolumn{4}{p{0.6\linewidth}}{
        \cellcolor[rgb]{0.561,0.737,0.561}{
            \makecell[{{p{\linewidth}}}]{
                \texttt{\tiny{[P1$\rangle$GM]}}
                \texttt{\{"current\_mode": "probing", "input\_rationale": "I have eight examples now: ({[}{-}1, 0, 1{]}, {-}1), ({[}2, 1{]}, 0), ({[}3, 2, 1{]}, 0), ({[}1, 1{]}, 0), ({[}{]}, 0), ({[}1, 2, 3{]}, 0), ({[}0{]}, 0), and ({[}{-}4{]}, {-}4). Let me analyze these. For {[}{-}1, 0, 1{]}, the result is {-}1. For {[}2, 1{]}, the result is 0. For {[}3, 2, 1{]}, the result is 0. For {[}1, 1{]}, the result is 0. For {[}{]}, the result is 0. For {[}1, 2, 3{]}, the result is 0. For {[}0{]}, the result is 0. For {[}{-}4{]}, the result is {-}4. This still doesn't immediately suggest a simple min/max or sum. Let me try another example to gather more data.", "hypothesis": "unknown", "confidence\_score": 80\}} \\
            }
        }
    }
    & & \\ \\

    \theutterance \stepcounter{utterance}  
    & & & \multicolumn{4}{p{0.6\linewidth}}{
        \cellcolor[rgb]{0.9,0.9,0.9}{
            \makecell[{{p{\linewidth}}}]{
                \texttt{\tiny{[P1$\langle$GM]}}
                \texttt{OUTPUT: ({[}5, {-}5, 5{]}, {-}5)} \\
            }
        }
    }
     \\ \\

    \theutterance \stepcounter{utterance}  
    & \multicolumn{4}{p{0.6\linewidth}}{
        \cellcolor[rgb]{0.561,0.737,0.561}{
            \makecell[{{p{\linewidth}}}]{
                \texttt{\tiny{[P1$\rangle$GM]}}
                \texttt{\{"current\_mode": "confirming", "input\_rationale": "I have nine examples now. Let's look at the pattern: \textbackslash{}n1. {[}{-}1, 0, 1{]} {-}> {-}1 (min is {-}1)\textbackslash{}n2. {[}2, 1{]} {-}> 0 (min is 1, but result is 0? No, wait. Min is 1. Result is 0. This contradicts min.)\textbackslash{}n3. {[}3, 2, 1{]} {-}> 0 (min is 1. Result is 0. Contradicts min.)\textbackslash{}n4. {[}1, 1{]} {-}> 0 (min is 1. Result is 0. Contradicts min.)\textbackslash{}n5. {[}{]} {-}> 0\textbackslash{}n6. {[}1, 2, 3{]} {-}> 0 (min is 1. Result is 0. Contradicts min.)\textbackslash{}n7. {[}0{]} {-}> 0 (min is 0. Result is 0. Consistent with min.)\textbackslash{}n8. {[}{-}4{]} {-}> {-}4 (min is {-}4. Result is {-}4. Consistent with min.)\textbackslash{}n9. {[}5, {-}5, 5{]} {-}> {-}5 (min is {-}5. Result is {-}5. Consistent with min.)\textbackslash{}n\textbackslash{}nWait, let's re{-}examine. \textbackslash{}nFor {[}2, 1{]}, min is 1} \\
            }
        }
    }
    & & \\ \\

    \theutterance \stepcounter{utterance}  
    & & & \multicolumn{4}{p{0.6\linewidth}}{
        \cellcolor[rgb]{0.9,0.9,0.9}{
            \makecell[{{p{\linewidth}}}]{
                \texttt{\tiny{[P1$\langle$GM]}}
                \texttt{Game Violation: ParseError: Response is not valid JSON.} \\
\texttt{Output ONLY:} \\
\texttt{One JSON with confidence\_score, current\_mode, input\_rationale, and hypothesis keys ...} \\
\texttt{OR} \\
\texttt{SOLVE: ```python ... ```} \\
            }
        }
    }
     \\ \\

    \theutterance \stepcounter{utterance}  
    & \multicolumn{4}{p{0.6\linewidth}}{
        \cellcolor[rgb]{0.561,0.737,0.561}{
            \makecell[{{p{\linewidth}}}]{
                \texttt{\tiny{[P1$\rangle$GM]}}
                \texttt{\{"current\_mode": "probing", "input\_rationale": "I have nine examples. Let's look at the pattern more carefully. \textbackslash{}n1. {[}{-}1, 0, 1{]} {-}> {-}1. Min is {-}1. Sum is 0. \textbackslash{}n2. {[}2, 1{]} {-}> 0. Min is 1. Sum is 3. \textbackslash{}n3. {[}3, 2, 1{]} {-}> 0. Min is 1. Sum is 6. \textbackslash{}n4. {[}1, 1{]} {-}> 0. Min is 1. Sum is 2. \textbackslash{}n5. {[}{]} {-}> 0. \textbackslash{}n6. {[}1, 2, 3{]} {-}> 0. Min is 1. Sum is 6. \textbackslash{}n7. {[}0{]} {-}> 0. Min is 0. Sum is 0. \textbackslash{}n8. {[}{-}4{]} {-}> {-}4. Min is {-}4. Sum is {-}4. \textbackslash{}n9. {[}5, {-}5, 5{]} {-}> {-}5. Min is {-}5. Sum is 5. \textbackslash{}n\textbackslash{}nIt seems like the function returns the minimum value if the minimum is negative, and 0 otherwise? Let's check:\textbackslash{}n1. min({-}1, 0, 1) = {-}1. Negative. Returns {-}1. Correct.\textbackslash{}n2. min(2, 1) = 1} \\
            }
        }
    }
    & & \\ \\

    \theutterance \stepcounter{utterance}  
    & & & \multicolumn{4}{p{0.6\linewidth}}{
        \cellcolor[rgb]{0.9,0.9,0.9}{
            \makecell[{{p{\linewidth}}}]{
                \texttt{\tiny{[P1$\langle$GM]}}
                \texttt{Game Violation: ParseError: Response is not valid JSON.} \\
\texttt{Output ONLY:} \\
\texttt{One JSON with confidence\_score, current\_mode, input\_rationale, and hypothesis keys ...} \\
\texttt{OR} \\
\texttt{SOLVE: ```python ... ```} \\
            }
        }
    }
     \\ \\

    \theutterance \stepcounter{utterance}  
    & \multicolumn{4}{p{0.6\linewidth}}{
        \cellcolor[rgb]{0.561,0.737,0.561}{
            \makecell[{{p{\linewidth}}}]{
                \texttt{\tiny{[P1$\rangle$GM]}}
                \texttt{\{"current\_mode": "confirming", "input\_rationale": "Based on the examples: {[}{-}1, 0, 1{]} {-}> {-}1 (min is {-}1), {[}2, 1{]} {-}> 0 (min is 1, but result is 0), {[}3, 2, 1{]} {-}> 0 (min is 1, result 0), {[}1, 1{]} {-}> 0 (min is 1, result 0), {[}{]} {-}> 0, {[}1, 2, 3{]} {-}> 0 (min is 1, result 0), {[}0{]} {-}> 0, {[}{-}4{]} {-}> {-}4 (min is {-}4), {[}5, {-}5, 5{]} {-}> {-}5 (min is {-}5). The pattern appears to be: return the minimum value if it is negative, otherwise return 0. This is equivalent to max(0, min(items)) if we consider 0 as the floor for non{-}negative results, but actually it's min(0, min(items)) if min is negative, else 0. Wait, for {[}2,1{]}, min is 1. Result is 0. For {[}{-}1,0,1{]}, min is {-}1. Result is {-}1. For {[}{-}4{]}, min is {-}4. Result is {-}4. For {[}5,{-}5,5{]}, min is {-}5. Result is {-}5. For {[}0{]}, min} \\
            }
        }
    }
    & & \\ \\

    \theutterance \stepcounter{utterance}  
    & & & \multicolumn{4}{p{0.6\linewidth}}{
        \cellcolor[rgb]{0.9,0.9,0.9}{
            \makecell[{{p{\linewidth}}}]{
                \texttt{\tiny{[P1$\langle$GM]}}
                \texttt{Game Violation: ParseError: Response is not valid JSON.} \\
\texttt{Output ONLY:} \\
\texttt{One JSON with confidence\_score, current\_mode, input\_rationale, and hypothesis keys ...} \\
\texttt{OR} \\
\texttt{SOLVE: ```python ... ```} \\
            }
        }
    }
     \\ \\

    \theutterance \stepcounter{utterance}  
    & \multicolumn{4}{p{0.6\linewidth}}{
        \cellcolor[rgb]{0.561,0.737,0.561}{
            \makecell[{{p{\linewidth}}}]{
                \texttt{\tiny{[P1$\rangle$GM]}}
                \texttt{\{"current\_mode": "confirming", "input\_rationale": "The pattern is clear: return the minimum value if it is negative, otherwise return 0. This can be expressed as min(0, min(items)) if items is not empty, else 0. Or simply: if the minimum of the list is negative, return it; otherwise return 0. This is equivalent to min(0, min(items)) for non{-}empty lists, and 0 for empty lists.", "hypothesis": "def f(items):\textbackslash{}n    \textbackslash{}"\textbackslash{}"\textbackslash{}"Returns the minimum value if it is negative, otherwise returns 0.\textbackslash{}"\textbackslash{}"\textbackslash{}"\textbackslash{}n    if not items:\textbackslash{}n        return 0\textbackslash{}n    return min(0, min(items))", "confidence\_score": 95\}} \\
            }
        }
    }
    & & \\ \\

    \theutterance \stepcounter{utterance}  
    & & & \multicolumn{4}{p{0.6\linewidth}}{
        \cellcolor[rgb]{0.9,0.9,0.9}{
            \makecell[{{p{\linewidth}}}]{
                \texttt{\tiny{[P1$\langle$GM]}}
                \texttt{OUTPUT: ({[}{-}1{]}, {-}1)} \\
            }
        }
    }
     \\ \\

    \theutterance \stepcounter{utterance}  
    & \multicolumn{4}{p{0.6\linewidth}}{
        \cellcolor[rgb]{0.561,0.737,0.561}{
            \makecell[{{p{\linewidth}}}]{
                \texttt{\tiny{[P1$\rangle$GM]}}
                \texttt{\{"current\_mode": "confirming", "input\_rationale": "The new example ({[}{-}1{]}, {-}1) is consistent with the hypothesis: min({[}{-}1{]}) is {-}1, which is negative, so return {-}1. All previous examples are also consistent. I am confident in this solution.", "hypothesis": "def f(items):\textbackslash{}n    \textbackslash{}"\textbackslash{}"\textbackslash{}"Returns the minimum value if it is negative, otherwise returns 0.\textbackslash{}"\textbackslash{}"\textbackslash{}"\textbackslash{}n    if not items:\textbackslash{}n        return 0\textbackslash{}n    return min(0, min(items))", "confidence\_score": 100\}} \\
            }
        }
    }
    & & \\ \\

    \theutterance \stepcounter{utterance}  
    & & & \multicolumn{4}{p{0.6\linewidth}}{
        \cellcolor[rgb]{0.9,0.9,0.9}{
            \makecell[{{p{\linewidth}}}]{
                \texttt{\tiny{[P1$\langle$GM]}}
                \texttt{OUTPUT: ({[}{-}1, 10, 26, 1, {-}1{]}, {-}2)} \\
            }
        }
    }
     \\ \\

    \theutterance \stepcounter{utterance}  
    & \multicolumn{4}{p{0.6\linewidth}}{
        \cellcolor[rgb]{0.561,0.737,0.561}{
            \makecell[{{p{\linewidth}}}]{
                \texttt{\tiny{[P1$\rangle$GM]}}
                \texttt{\{"current\_mode": "probing", "input\_rationale": "The new example ({[}{-}1, 10, 26, 1, {-}1{]}, {-}2) contradicts my previous hypothesis. My hypothesis predicted {-}1 (the minimum), but the actual output is {-}2. This means the function is not simply returning the minimum. Let's re{-}evaluate. \textbackslash{}nPrevious examples:\textbackslash{}n{[}{-}1, 0, 1{]} {-}> {-}1\textbackslash{}n{[}2, 1{]} {-}> 0\textbackslash{}n{[}3, 2, 1{]} {-}> 0\textbackslash{}n{[}1, 1{]} {-}> 0\textbackslash{}n{[}{]} {-}> 0\textbackslash{}n{[}1, 2, 3{]} {-}> 0\textbackslash{}n{[}0{]} {-}> 0\textbackslash{}n{[}{-}4{]} {-}> {-}4\textbackslash{}n{[}5, {-}5, 5{]} {-}> {-}5\textbackslash{}n{[}{-}1{]} {-}> {-}1\textbackslash{}n{[}{-}1, 10, 26, 1, {-}1{]} {-}> {-}2\textbackslash{}n\textbackslash{}nLet's look at the sum of elements for the new example: {-}1 + 10 + 26 + 1 {-} 1 = 35. Not 2.\textbackslash{}nLet's look at the count of negative numbers: 2. Not 2.\textbackslash{}nLet's look at the number of elements: 5. Not 2.\textbackslash{}nLet's look at the minimum: {-}1. Not 2.\textbackslash{}nLet's look at the second minimum} \\
            }
        }
    }
    & & \\ \\

    \theutterance \stepcounter{utterance}  
    & & & \multicolumn{2}{p{0.3\linewidth}}{
        \cellcolor[rgb]{0.9,0.9,0.9}{
            \makecell[{{p{\linewidth}}}]{
                \texttt{\tiny{[GM$|$GM]}}
                \texttt{Game Verdict: LOSS (turn limit reached)} \\
            }
        }
    }
    & & \\ \\

\end{supertabular}
}

\end{document}
